\subsection{Таймеры}
\begin{itemize}
  \item \verb|Timer.new(period, fn)| — создаёт \hyperlink{term:timer}{периодический таймер}; \verb|period| — \hyperlink{term:period}{период} в секундах; \verb|fn| вызывается каждый \hyperlink{term:tick}{тик}.
  \item \verb|timer:start()|, \verb|timer:stop()| — запуск/остановка таймера.
  \item \verb|Timer.callLater(delay, fn)| — единичный вызов \verb|fn| через \verb|delay| секунд.
  \item Таймеры и \hyperlink{term:state}{состояния} делайте \hyperlink{term:scope}{глобальными} (без \verb|local|), иначе сборщик мусора их удалит.
\end{itemize}
\textbf{Пример периодического таймера (мигание):}
\begin{codefile}{lua/examples/05_leds/leds_timers_blink.lua}
\end{codefile}
\textbf{Пояснение.} В примере функция \verb|updateBlink()| инвертирует флаг \verb|on| и по таймеру \verb|blinkTimer| раз в 0{,}5 секунды включает или выключает все светодиоды. Основной принцип — периодически вызывать одну и ту же функцию, чтобы реализовать мигание без блокирующих задержек.
\textbf{Пример отложенного вызова (через 3 секунды):}
\begin{codefile}{lua/examples/05_leds/leds_timers_delayed.lua}
\end{codefile}
\textbf{Пояснение.} Скрипт создаёт объект \verb|Ledbar| и с помощью \verb|Timer.callLater| один раз через 3 секунды включает синий цвет на всех светодиодах. Основной принцип — выполнить действие с задержкой, не блокируя основной поток выполнения.
\textbf{Отложенный таймер внутри обработчика событий:}
\begin{codefile}{lua/examples/05_leds/leds_timers_callback_delayed.lua}
\end{codefile}
\textbf{Пояснение.} Здесь функция \verb|callback(event)| реагирует на событие \verb|Ev.LOW_VOLTAGE2| и внутри неё запускает \verb|Timer.callLater|, который через секунду заливает линейку красным цветом. Основной принцип — сочетать обработку событий и отложенные таймеры, оставаясь в неблокирующей модели.
